\section*{Apêndice A --- Visualizações históricas e leitura dos resultados}
\label{ap:visualizacoes}

Os gráficos de trajetória usam o ano do snapshot no eixo horizontal e uma medida CK por painel. São apresentadas separadamente a média anual entre projetos e as linhas individuais dos projetos criados entre 2009 e 2013. O ponto parcial de 2026 é identificado; lacunas de cobertura não são substituídas por zero. Mudanças na média anual podem decorrer da composição de projetos disponíveis, enquanto a trajetória individual acompanha o mesmo repositório.

As tabelas por época apresentam uma contribuição por repositório, com quantidades de projetos e snapshots. O total de projetos únicos não deve ser obtido pela soma das épocas, pois um mesmo projeto pode atravessá-las. Os contrastes por evidência distinguem presença observada, ausência de evidência e transição dentro da época.

Nas tabelas antes/depois, \texttt{n\_a} é o número de projetos com medição anterior e \texttt{n\_b} o número com medição posterior. Os pares correspondem aos mesmos repositórios. A diferença é calculada dentro de cada par antes de obter a média; para os mesmos pares e pesos, ela corresponde à média posterior menos a média anterior. Os SHA e as datas permitem conferir cada versão utilizada.

As sete medidas conservam suas unidades originais. Não se cria uma nota composta de maturidade. Não se apresentam p-valores, deltas de Cliff ou afirmações automáticas de melhoria. Arquivos CSV, gráficos PNG/SVG, referências fixadas, evidências, diagnósticos e hashes do instrumento integram o pacote de reprodução. Os relatórios dos pilotos transversais anteriores permanecem identificados separadamente.

\section*{Apêndice B --- Cobertura do piloto histórico}

O piloto compreendeu os primeiros dez candidatos sorteados de cada coorte. Dos 720 cortes anuais, 373 não tinham histórico disponível até o corte, 346 tiveram extração CK concluída e um não continha fontes Java. Os 347 snapshots com árvore disponível corresponderam a 127 extrações distintas, em razão da reutilização de fontes idênticas. A tabela conserva os resultados do piloto corrigido, anterior à ampliação.

\begin{table}[htbp]
\centering
\caption{Cobertura do piloto histórico por coorte de criação.}
\label{tab:piloto-historico-cobertura}
\small
\begin{tabular}{lrrrr}
\hline
Coorte & Candidatos & Snapshots disponíveis & Elegíveis & Projetos elegíveis \\
\hline
E1 & 10 & 155 & 77 & 5 \\
E2 & 10 & 101 & 31 & 3 \\
E3 & 10 & 75 & 8 & 1 \\
E4 & 10 & 16 & 3 & 3 \\
\hline
Total & 40 & 347 & 119 & 12 \\
\hline
\end{tabular}
\end{table}

A reexecução utilizou quatro trabalhadores e consumiu aproximadamente 355 segundos, reutilizando cerca de 505 MiB de clones e produzindo aproximadamente 44 MiB de extrações. Esses valores descrevem o ambiente e o cache daquela execução, sem estimar o custo de uma coleta inteiramente nova.
